\section{A partition-based second step}
\label{sec:partition-step}

The new construction keeps the first spatial head and appends another
interval of length \(h=\tfrac14\log(8/7)\). The resulting partition has lengths
\[
 (\ell_0,\ell_1,\ell_2)=(\log7,h,h),\qquad
 L_2=\tfrac12\log56=2.0126758453675746166785\ldots,
\]
with retained dimensions \((256,32,32)\). Thus the old verification space at
this step consists of the preceding 288 spatial modes. It is not replaced
by a new global polynomial space.

\subsection{Complete outputs and inherited data}
On each output interval, the truncated gamma output of every retained input
has the form
\[
 p(v)+\sum_c p_c(v)\log|v-c|,\qquad 0<v<1,
\]
where every logarithmic centre lies outside the open interval. Arithmetic
outputs are polynomials restricted to their exact translated windows.
The endpoints are rational combinations of prime logarithms; their equality
is decided symbolically before interval division. At this horizon the
generator active set is \(\{2,3,4,5,7\}\).

Products of these outputs are integrated without an output projection.
Polynomial--logarithm moments follow by polynomial division and integration
by parts. For two distinct centres, the remaining integral
\[
 \int_0^1\frac{\log|v-d|}{v-c}\,dv
\]
has the primitive
\[
 \log|d-c|\log|v-c|
 -\Rea\operatorname{Li}_2\!\left(\frac{v-c}{d-c}\right).
\]
Endpoint limits are taken before evaluating this expression, and divergent
terms are omitted only when their coefficients vanish identically.
For \(c=d\notin[0,1]\), the primitive is
\(\tfrac12\log^2|v-c|\). The same-endpoint square-log cases follow directly
from integration by parts. All sign computations use Arb enclosures;
quadrature is used only for consistency checks. The expanded checks cover
52 endpoint and exterior-log moment values, including interval-valued centres
near $-1$ and $2$, where one endpoint logarithm encloses zero. Squares are
evaluated by multiplication rather than a general power operation. The new
LDL validator also checks finiteness before recording a failed pivot.
Three-interval squared-output Gram checks agree with independent quadrature
to better than $7\times10^{-77}$ on the selected small examples. These checks
support implementation consistency, not a separate audit of the full build.

The first-step head and its old/old output Gram blocks are inherited from
the hash-verified R7 archive. The new cross blocks and new output interval
are reconstructed. This is an incremental construction, not an independent
reconstruction of the entire first step. The inherited exterior-log series
and the new exact-log representation need not give identical model outputs,
so their mismatch is included in the new Gram error.

More precisely, let \(d_s\) bound the seed output error, let
\(t_s\) be the trace of its full model Gram, and let \(\eta\) bound the
new exact-profile model error. Each of the two inherited output restrictions
differs from the new model by at most \(d_s+\eta\). Therefore the sum of the
inherited Gram block errors is bounded by
\[
 e_s=2\{2(d_s+\eta)\sqrt{t_s}+(d_s+\eta)^2\}.
\]
Writing \(\widetilde G\) for the assembled full Gram and \(n=320\), a valid
bound for the norm of the new model output is
\[
 b_{\rm out}=\sqrt{\operatorname{tr}\widetilde G+n e_s}.
\]
Consequently the full Gram and leakage errors are at most
\[
 e_G=e_s+2\eta b_{\rm out}+\eta^2,\qquad
 e_E=e_G+2\eta\|\widetilde H\|_F+\eta^2.
\]
The same bounds cover a restriction to any union of output intervals.
These inequalities also explain why blindly combining differently
approximated inherited and new blocks would be insufficient.

\subsection{The multi-interval complement}
Let \(K_i\) omit the first \(N_i\) local Legendre modes. The diagonal gamma
tail bound is the corrected \(g_{\ell_i,N_i}\). For two intervals separated
by a positive gap \(d_{ij}\), the singular gamma kernel has the
Hilbert--Schmidt bound
\[
 b_{ij}^{\rm sing}\le
 \frac12\sqrt{\log
 \frac{(d_{ij}+\ell_i)(d_{ij}+\ell_j)}
 {d_{ij}(d_{ij}+\ell_i+\ell_j)}}.
\]
For adjacent intervals use \(\pi/2\); the smaller bound may be used when
both apply. Integration by parts on either tail gives
\[
 b_{ij}^{\rm sm}\le
 \frac{\sqrt{\ell_i\ell_j}}4
 \min\left\{\frac{\ell_i}{\sqrt{N_i(N_i+1)}},
             \frac{\ell_j}{\sqrt{N_j(N_j+1)}}\right\}
 \sup_{[0,L_2]}|w_0'|.
\]
Form the symmetric comparison matrix with diagonal \(g_{\ell_i,N_i}\)
and off-diagonal entries \(-b_{ij}^{\rm sing}-b_{ij}^{\rm sm}\).
Subtract a newly verified continuum arithmetic Schur bound \(\rho_2 I\).
A finite interval LDL test then bounds its least eigenvalue from below,
which bounds the full three-interval complement. At the recorded parameters
this lower bound is greater than \(0.4109027593994060\).

The complete-output Schur test, the old-energy comparison, and the enlarged
graph residual test use the same reductions as for the first step.
For the latter two tests the old output is the union of the first two
intervals. A rational Galerkin proposal is extended by zero on their omitted
modes. Every translated old source window is tested for overlap; the
arithmetic cross bound uses the disjoint-window improvement only when this
test passes, and otherwise uses the sum of the delay norms.

\subsection{The resulting certificate boundary}
\label{sec:second-result}
The new rational continuation has 288 old input columns and 32 new output
rows, rounded to 80 decimal digits and then treated as exact rationals.
It is extended by zero on the two old polynomial complements. The new
complete-output archive has
\[
 \eta<5.494\times10^{-56},\qquad
 e_G<2.376\times10^{-53},\qquad e_E<3.202\times10^{-53}.
\]
Its corrected construction took approximately 764 seconds in the recorded
environment. The compressed archive occupies 297,151,691 bytes and is
stored outside Git.

\begin{proposition}[Second-step old-energy comparison]
Let \(A_2=Q_{0,L_q}\), let \(F_2=Q^\gamma_{0,h}\), and let \(B_2\) be their
complete spatial cross block at \(L_2\). For the new rational \(J_2\), put
\[
 H_{J_2}=A_2+B_2^*J_2+J_2^*B_2+J_2^*F_2J_2.
\]
In the working normalization, on the full old form domain,
\[
 H_{J_2}\succeq10^{-8}A_2.
\]
\end{proposition}
\begin{proof}[Reduction and numerical premises]
Use the comparison reduction of Proposition~\ref{prop:old-energy-comparison}
with old head dimension 288 and old output equal to the union of the first
two intervals. The old-tail lower bound for
\(H_{J_2}-10^{-8}A_2\) exceeds \(0.8299288982523047\).
The guarded error is below \(9.270\times10^{-51}\).
All 288 pivots pass at both 2,048 and 4,096 bits, with minimum pivot lower
bound approximately \(2.2736098\times10^{-5}\). As before, that pivot
diagnostic is not a spectral floor. The same sufficient test at
\(\mu=10^{-7}\) fails; this does not disprove that stronger inequality.
\end{proof}

The remaining residual premise has not been certified. At \(\theta=0.9\),
\(0.95\), and \(0.99\), the scalar-complement graph tests fail at zero-based
pivots 55, 155, and 174, respectively. The absolute sufficient test requesting
\(Q_{0,L_2}\succeq10^{-36}I\) fails at pivot 182. These outcomes are reproduced
at 2,048 and 4,096 bits; they are not numerical witnesses against either the
operator or the desired inequalities.

The separate output Grams allow a further analytic refinement. Let
\(\Gamma_\theta\) be the three-by-three comparison matrix for the enlarged
graph form, including its arithmetic bound. For any positive rational
weights \(w_i\), define certified lower bounds
\[
 \beta_i\le(\Gamma_\theta)_{ii}
       -\sum_{j\ne i}|(\Gamma_\theta)_{ij}|\,\frac{w_j}{w_i}.
\]
The inequality \(2ab\le(w_j/w_i)a^2+(w_i/w_j)b^2\) proves a separate
tail floor \(\beta_i\) on each interval. Assume all \(\beta_i>0\).
Let \(\widetilde E_i\) be the model leakage Gram into interval \(i\)'s
complement before the graph change, and let \(T\) be the finite graph map.
Put \(D_i=\operatorname{diag}(I_o,\lambda I_n)\) for either old interval
and \(D_2=\operatorname{diag}(\lambda I_o,I_n)\) for the new one,
where \(\lambda=\theta^{-1/2}\).
The sufficient finite test becomes
\[
 \widetilde M_\theta-
 \sum_i\frac{D_iT^*\widetilde E_iTD_i}{\beta_i}
 -
 \left(2\lambda\eta\|T\|_F^2+
       e_E\sum_i\frac{\|TD_i\|_F^2}{\beta_i}\right)I
 \succ0.
\]
This follows by eliminating the separately bounded tails. It retains more
of the interval information than a single common floor.

Ten proposed weight triples were tested at each of \(\theta=0.99\) and
\(0.9999\), at 1,024 bits. None passed.
\begin{samepage}
\noindent For example, at \(\theta=0.9999\)
the exact rational triple
\[
 (1,\;0.6335965029937086,\;0.3895577505597481)
\]
gives tail floors greater than \(0.5318716111\), \(0.2621540228\), and
\(0.0963354365\), respectively, with error budget below
\(1.654\times10^{-49}\). The guarded pivot at index 286 is negative,
approximately \(-0.9413506220\). A separate 2,048-bit replay uses this
same triple. This finite search neither excludes other weights nor
proves failure of the residual inequality.
\end{samepage}

Thus the first-step induction premises are complete, while at the second
step only the old-energy premise has been established. No new full-operator
floor or positive-shift interval is claimed at \(L_2\). The largest certified
horizon in the current package remains the separate \(1.98\) result R6.
The next concrete experiment is to enlarge the verification spaces on the
short slabs and strengthen their joint complement estimate, reusing the
stored complete outputs wherever possible. Pivot locations are diagnostics
for that experiment, not a localization theorem for an obstruction.

\subsection{What is needed beyond the radius-three implementation}
The output and tail routines permit general finite partitions. The current
driver appends the second quarter slab using one gamma Taylor profile about
zero, and therefore enforces \(L<3\). This is a numerical
representation restriction. It is not an all-depth induction theorem.
A local representation provides a concrete way to remove that restriction.

The central gamma off-diagonal kernel satisfies the exact identity
\[
 q(t)=2\cosh(t/2)
       -\frac{e^{-t/2}}{1-e^{-2t}}
     =2\cosh(t/2)-\sum_{k=0}^{\infty}e^{-(1/2+2k)t},
 \qquad t>0.
\]
On any interval the two exponential functions \(e^{x/2},e^{-x/2}\)
span the restriction of the finite-rank hyperbolic term. If each local
verification space includes these functions, that term has zero compression
and zero coupling on its orthogonal complement. This statement concerns
the exact enlarged span; replacing the exponentials by polynomial
approximations requires a further error bound.

For intervals with gap \(d>0\), after retaining \(K\) exponential terms the
remaining cross operator has norm at most
\[
 \sqrt{\ell_i\ell_j}\,
 \frac{e^{-(1/2+2K)d}}{1-e^{-2d}}.
\]
This follows by summing the positive geometric remainder pointwise and then
using its Hilbert--Schmidt bound. Each retained exponential is separable
across the two intervals. Local Taylor profiles can still handle each
diagonal and adjacent pair whose total length is less than three.
Thus a fixed short-slab partition together with the separated-interval
expansion has no global radius-three barrier.

This representation has not yet been implemented or certified here.
It also does not control the arithmetic contribution or prove the
old-energy and residual inequalities indefinitely. Those comparisons,
their required verification dimensions, and prevention of finite-depth
accumulation remain the substantive induction problem.
